\documentclass[11pt,a4paper]{article}
\usepackage[utf8]{inputenc}
\usepackage[spanish]{babel}
\usepackage{amsmath, amssymb}
\usepackage{geometry}
\geometry{top=2.5cm, bottom=2.5cm, left=2.5cm, right=2.5cm}

\begin{document}

\section*{Exámenes de práctica del capítulo 10}

\subsection*{Examen 1}

\begin{itemize}
    \item[\textbf{1.}] Un hombre grande se sienta en una silla de cuatro patas, con sus pies levantados del piso. La masa combinada del hombre y la silla es \(95.0\,\text{kg}\). Si las patas de la silla son circulares y tienen un radio de \(0.500\,\text{cm}\) en el fondo, ¿qué presión ejerce cada pata sobre el piso?
    \vspace*{6cm}
    
    \item[\textbf{2.}] Blaise Pascal duplicó el barómetro de Torricelli usando un vino tinto Bordeaux, de \(984\,\text{kg/m}^3\) de densidad, como el líquido de trabajo.
    \begin{enumerate}
        \item[\textbf{(a)}] ¿Cuál fue la altura \(h\) de vino para presión atmosférica normal?
        \item[\textbf{(b)}] ¿Esperaría que el vacío sobre la columna sea tan bueno como para el mercurio?
    \end{enumerate}
    
    % ==========================================
    % ESPACIO PARA LA FIGURA DEL EJERCICIO 2 (Barómetro de Pascal)
    % ==========================================
    % \begin{figure}[h]
    % \centering
    % \includegraphics[width=0.5\textwidth]{tu_imagen.png}
    % \caption{Barómetro con vino}
    % \end{figure}
    \vspace*{7cm}
\end{itemize}

\newpage

\begin{itemize}
    \item[\textbf{3.}] Sobre un dique de altura \(h\) cae agua con un caudal de masa \(R\), en unidades de kilogramos por segundo.
    \begin{enumerate}
        \item[\textbf{(a)}] Demuestre que la potencia disponible a causa del agua es
        \[
        P = Rgh
        \]
        donde \(g\) es la aceleración en caída libre.
        \item[\textbf{(b)}] Cada unidad hidroeléctrica en el dique Grand Coulee toma agua en una cantidad de \(8.50 \times 10^5\,\text{kg/s}\) desde una altura de \(87.0\,\text{m}\). La potencia desarrollada por la caída de agua se convierte en energía eléctrica con una eficiencia del \(85.0\%\). ¿Cuánta energía eléctrica produce cada unidad hidroeléctrica?
    \end{enumerate}
    \vspace*{7cm}
    
    \item[\textbf{4.}] Evangelista Torricelli fue la primera persona en darse cuenta de que los seres humanos viven en el fondo de un océano de aire. Él conjeturó correctamente que la presión de la atmósfera es atribuible al peso del aire. La densidad del aire a \(0^\circ\text{C}\) en la superficie de la Tierra es \(1.29\,\text{kg/m}^3\). La densidad disminuye al aumentar la altitud (a medida que la atmósfera se adelgaza). Por otra parte, si se supone que la densidad es constante a \(1.29\,\text{kg/m}^3\) hasta cierta altura \(h\) y es cero sobre dicha altura, en tal caso \(h\) representaría la profundidad del océano de aire.
    \begin{enumerate}
        \item[\textbf{(a)}] Use este modelo para determinar el valor de \(h\) que da una presión de \(1.00\,\text{atm}\) en la superficie de la Tierra.
        \item[\textbf{(b)}] ¿La cima del monte Everest se elevaría sobre la superficie de tal atmósfera?
    \end{enumerate}
    \vspace*{7cm}
\end{itemize}

\newpage

\subsection*{Examen 2}

\begin{itemize}
    \item[\textbf{1.}] Estime la masa total de la atmósfera de la Tierra. (El radio de la Tierra es \(6.37 \times 10^6\,\text{m}\), y la presión atmosférica en la superficie es \(1.013 \times 10^5\,\text{Pa}\).)
    \vspace*{6cm}
    
    \item[\textbf{2.}] La fuerza gravitacional que se ejerce sobre un objeto sólido es \(5.00\,\text{N}\). Cuando el objeto se suspende de una balanza de resorte y se sumerge en agua, la lectura en la balanza es \(3.50\,\text{N}\). Encuentre la densidad del objeto.
    
    % ==========================================
    % ESPACIO PARA LA FIGURA DEL EJERCICIO 2 (Balanzas)
    % ==========================================
    % \begin{figure}[h]
    % \centering
    % \includegraphics[width=0.5\textwidth]{tu_imagen_balanza.png}
    % \caption{Objeto sumergido}
    % \end{figure}
    \vspace*{7cm}
\end{itemize}

\newpage

\begin{itemize}
    \item[\textbf{3.}] Se ha colocado un recubrimiento fino de glicerina de espesor \(1.50\,\text{mm}\) entre dos portaobjetos de microscopio de \(1.00\,\text{cm}\) de ancho y \(4.00\,\text{cm}\) de largo. Encuentre la fuerza requerida para jalar de uno de los portaobjetos del microscopio a una velocidad constante de \(0.300\,\text{m/s}\) en relación con el otro portaobjetos.
    \vspace*{6cm}
    
    \item[\textbf{4.}] En 1983, Estados Unidos comenzó a acuñar la moneda de un centavo a partir de zinc revestido de cobre en lugar de cobre puro. La masa del antiguo penique de cobre es \(3.083\,\text{g}\) y el del nuevo centavo es de \(2.517\,\text{g}\). Considere que la densidad del cobre es \(8.960\,\text{g/cm}^3\) y la del zinc es \(7.133\,\text{g/cm}^3\). Las monedas nueva y antigua tienen el mismo volumen. Calcule el porcentaje de zinc (por volumen) en el nuevo centavo.
    \vspace*{7cm}
\end{itemize}

\end{document}
